\chapter{Compound types and collections}
\label{chap:collections}

Every variable of the previous chapters held one value: a number, a character,
a string. A program that keeps the grades of a class of thirty students would
need thirty variables, and a program that reads numbers until the user stops
cannot know how many variables to declare. Real programs handle data in bulk,
and the data comes in groups.

A \textit{compound type} groups several values into one. This chapter presents
the groups that Ymir provides without any declaration: \textit{tuples}, which
bind a few values of any types together; \textit{arrays} and \textit{slices},
which hold a sequence of values of the same type; \textit{maps}, which look up
a value by a key; and \textit{options}, which hold a value or nothing at all.
Ranges, met in Chapter~\ref{chap:control}, turn out to be values of this kind
too. A type that holds many values of the same kind, such as an array, a slice
or a map, is called a \textit{collection}, and \token{for} goes through each of
them as it goes through a range. The chapter ends with a program that keeps the
grades of a class.

A collection whose elements can change raises questions of its own: who else
sees the change, and what a copy is. This chapter gives the rule needed to
change a collection, and the next chapter explains it.

Part II covers the compound types exhaustively, in
Chapter~\ref{chap:compound}.

\minitoc%

\vfill\pagebreak
\section{Grouping values: tuples}
\label{sec:collections:tuples}

A function returns one value. Dividing 17 by 5 gives two answers, though: a
quotient of 3 and a remainder of 2. A \textit{tuple} groups a fixed number of
values into a single one, so that the two answers travel together: a function
can return them, a variable can hold them, and another function can receive
them.

\lstinputlisting[style=coloredverbatim, caption={A function that returns two values, \textit{examples/collections/divide.yr}}, label=lst:collections:divide]{examples/collections/divide.yr}
\vspace{-10pt}%
\begin{lstlisting}[style=bashVerb]
(3, 2)
quotient 3, remainder 2
100 = 7 * 14 + 2
\end{lstlisting}

A tuple is written as its values between parentheses, separated by commas, as
\token{(a / b, a \% b)} on line~4. Its type is written the same way, with a type
in place of each value: \token{(i32, i32)}, on line~3, is a pair of
\token{i32}. \token{println} prints a tuple the way it is written, on line~9.

\subsection{Reading the fields}

The values of a tuple are its \textit{fields}, numbered from 0 in the order they
are written. A dot followed by the number of a field reads it: on line~10,
\token{result.0} is the quotient and \token{result.1} the remainder. The number
of fields of a tuple is its \textit{arity}, and is part of its type. A pair has
no field 2, and asking for one is refused when compiling.

\begin{lstlisting}[style=coloredverbatimError, escapechar=@, caption=A field that does not exist]
use std::io;

fn main() {
  let pair = (3, 4);
  println(pair.0 + @\hb{pair.2}@); // error
}
\end{lstlisting}

The compiler says \errmsg{E4214}{tuple access out of bound(2), tuple arity is
  2}.

The fields of a tuple need not have the same type. In the next listing,
\token{person} is a \token{([c8], i32, f64)}: a string, an integer and a
floating point number, which stand for a name, an age and a height.

\lstinputlisting[style=coloredverbatim, caption={A tuple of three types, \textit{examples/collections/person.yr}}, label=lst:collections:person]{examples/collections/person.yr}
\vspace{-10pt}%
\begin{lstlisting}[style=bashVerb]
Alice, 31 years old, 1.68 m tall
Alice: 1.68 m
Alice
31
1.68
\end{lstlisting}

\subsection{Taking a tuple apart}

\token{result.0} and \token{result.1} say where a value is, not what it is. A
tuple can instead be taken apart into variables with names of their own. On
line~12 of Listing~\ref{lst:collections:divide},
\token{let (q, r) = divide(100, 7);} declares two variables at once, one per
field, in order: \token{q} holds the quotient and \token{r} the remainder. The
number of names must be the arity of the tuple.

A field that the program does not need takes the wildcard \token{\_} in place of
a name, as the iterator of a loop does (Section~\ref{sec:control:for}). On
line~7 of Listing~\ref{lst:collections:person}, the age is left out, and only
\token{name} and \token{height} are declared.

\subsection{Passing the fields as arguments: \texttt{expand}}
\label{sec:collections:expand}

A function that takes three parameters cannot receive a tuple of three fields:
the call gives it one argument, where it expects three. The fields could be
passed one by one, as \token{volume(size.0, size.1, size.2)} does on line~9 of
the next listing, but the keyword \token{expand} says it in one word.

\lstinputlisting[style=coloredverbatim, caption={Expanding a tuple, \textit{examples/collections/volume.yr}}, label=lst:collections:volume]{examples/collections/volume.yr}
\vspace{-10pt}%
\begin{lstlisting}[style=bashVerb]
24
24
300
\end{lstlisting}

\token{expand size}, on line~10, replaces the tuple with its fields, separated by
commas, as if they had been written one by one. The call therefore receives
three arguments. The fields need not be all the arguments: on line~13,
\token{expand floor} gives the first two, and \token{10} the third.

Without \token{expand}, \token{volume(floor, 10)} is refused with
\errmsg{E4226}{the call operator is not defined for \errhole{} and
  \errhole{}}, and the note under it says that 3 parameters are expected, where
2 are given.

\token{expand} is not limited to calls. Inside the parentheses of a tuple, it
also stands for the fields of another tuple, which builds a longer tuple from
shorter ones.

\begin{lstlisting}[style=coloredverbatim]
let date = (28, 9, 2026);
let time = (14, 30);
let moment = (expand date, expand time);
println(moment);
let dated = (expand date, "Monday");
println(dated);
\end{lstlisting}
\vspace{-10pt}%
\begin{lstlisting}[style=bashVerb]
(28, 9, 2026, 14, 30)
(28, 9, 2026, Monday)
\end{lstlisting}

On line~3, \token{(expand date, expand time)} is read as if it were
\token{(28, 9, 2026, 14, 30)}: the three fields of \token{date} followed by the
two of \token{time}, a tuple of five fields. On line~5, the fields of
\token{date} are followed by a string, and \token{dated} is an
\token{(i32, i32, i32, [c8])}: the fields keep their types, and the string adds
its own.

\subsection{Going through the fields}

\token{for} goes through the fields of a tuple, from the first to the last, as
on line~10 of Listing~\ref{lst:collections:person}. With two names, the first
holds the number of the field, and the second the field itself.

\begin{lstlisting}[style=coloredverbatim]
let row = ("Alice", 31, 1.68, true);
for i, field in row {
  println(i, ": ", field, ", of type ", typeof(field)::typeid);
}
\end{lstlisting}
\vspace{-10pt}%
\begin{lstlisting}[style=bashVerb]
0: Alice, of type [c8]
1: 31, of type i32
2: 1.68, of type f64
3: true, of type bool
\end{lstlisting}

\token{typeof(field)::typeid} is the name of the type of \token{field}, as a
string; a later chapter comes back to it. It shows what sets this loop apart
from the others: the iterator has the type of the current field, so
\token{field} is a \token{[c8]} in the first iteration, an \token{i32} in the
second, an \token{f64} in the third and a \token{bool} in the last.

The compiler can only do that by copying the block of the loop once for each
field, with a different type for \token{field} in each copy. The number of
iterations is therefore always the arity of the tuple, and the block must be
valid for every type of field. \token{println} is; \token{field + 1} is not,
since a string cannot be added to a number.

\begin{lstlisting}[style=coloredverbatimError, escapechar=@, caption=A block that is not valid for every field]
use std::io;

fn main() {
  let person = ("Alice", 31, 1.68);
  for field in person {
    println(@\hb{field + 1}@); // error
  }
}
\end{lstlisting}

The compiler says \errmsg{E4224}{undefined operator + for types [c8] and i32}:
the copy of the block made for the first field, a string, does not compile,
although the copies for the two numbers would.

\subsection{Comparing tuples}

Two tuples of the same type are equal when their fields are equal, one by one:
\token{divide(17, 5) == (3, 2)} is \token{true}. \token{!=} is its opposite.
Tuples have no order, so \token{<} and \token{>} are not defined on them: there
is no natural way to say which of \token{(3, 2)} and \token{(2, 3)} comes first.

\subsection{A tuple of one value}

Parentheses also group operations, as in \token{(a + b) * 2}, so \token{(42)}
is the number 42, not a tuple. A tuple with a single field is written with a
comma after its value, \token{(42,)}, and its type likewise, \token{(i32,)}.
Such tuples are common in generic code, which the course comes to in later
chapters: a function written for tuples of any arity must work for an arity of
one as well, and a list of arguments collected into a tuple may hold a single
value. They also appear by mistake. When a compiler message names a type such as
\token{(i32,)} where a number was expected, the cause is often a stray comma
after a single value.

\vfill\pagebreak
\section{A fixed number of values: arrays}
\label{sec:collections:arrays}

The calendar of Section~\ref{sec:functions:calendar} named the months with a
\token{match} of twelve arms, and counted their days with another. Both are
tables: twelve values, looked up by the number of a month. An \textit{array} is
such a table. It holds a fixed number of values of the same type, its
\textit{elements}, stored one after the other and numbered from 0.

\lstinputlisting[style=coloredverbatim, caption={The months in two arrays, \textit{examples/collections/months.yr}}, label=lst:collections:months]{examples/collections/months.yr}
\vspace{-10pt}%
\begin{lstlisting}[style=bashVerb, escapechar=@+]
@+\textcolor{teal!80}{alice@dev:\mtilde}@+$ ./months
Month: 2
February has 28 days.
\end{lstlisting}

An array is written as its elements between square brackets, separated by
commas, as on lines~4 and~7. A long one can be spread over several lines. Its
type gives the type of the elements and their number, separated by a semicolon:
\token{days} is an \token{[i32 ; 12]}, and \token{names} an
\token{[[c8] ; 12]}, twelve strings. The number is part of the type. An array of
12 elements and one of 13 have different types, and an array never grows or
shrinks. February has 28 days here; the leap years of the calendar are left
out, to keep the listing short.

\subsection{Indices}

\token{days[i]} is the element of \token{days} at \textit{index} \token{i}. The
first element is at index 0, so the last of twelve is at index 11, as
Figure~\ref{fig:collections:array_indices} shows. The user counts the months
from 1, and the array from 0, hence the \token{month - 1} of line~10. The index
can be any integer expression.

\begin{figure}[h]
  \centering
  \begin{adjustbox}{max width=\linewidth}
    \begin{tikzpicture}[cell/.style={rectangle, draw=black!60, fill=background!40,
          minimum width=10mm, minimum height=7mm, font=\ttfamily}]
      \foreach \value [count=\i from 0] in {31, 28, 31, 30, 31, 30, 31, 31, 30, 31, 30, 31} {
        \node[cell] (c\i) at (\i * 1.0, 0) {\value};
        \node[cflab, below=1mm of c\i] {\i};
      }
      \node[cflab, left=3mm of c0] {\token{days}};
      \node[cflab, above=5mm of c1] (feb) {\token{days[1]}};
      \node[cflab, above=5mm of c11] (last) {\token{days[11]}};
      \draw[cfflow] (feb) -- (c1);
      \draw[cfflow] (last) -- (c11);
      \node[cflab, below=6mm of c5.south east] {indices, from 0 to \token{days.len - 1}};
    \end{tikzpicture}
  \end{adjustbox}
  \caption{\label{fig:collections:array_indices} The array \token{days} of
    Listing~\ref{lst:collections:months}. Its twelve elements are stored one
    after the other, and numbered from 0: the days of February are
    \token{days[1]}, those of December \token{days[11]}.}
\end{figure}


The number of elements of an array is its \textit{length}, \token{days.len}. It
is a \token{usize}, the unsigned integer type whose width is that of the
machine (Section~\ref{sec:types:integers}), since a length is never negative.

\subsection{Out of bounds}

\token{days} has no element at index 12. When the index is known while
compiling, the compiler checks it.

\begin{lstlisting}[style=coloredverbatimError, escapechar=@, caption=An index past the end of an array]
use std::io;

fn main() {
  let days = [31, 28, 31, 30, 31, 30, 31, 31, 30, 31, 30, 31];
  println("December has ", @\hb{days[12]}@, " days."); // error
}
\end{lstlisting}

The compiler says \errmsg{E4005}{slice access, index overflow 12us (len) <= 12
  (index)}: the length is 12, and a valid index must be below it.

On line~10 of Listing~\ref{lst:collections:months}, the index depends on what the
user types, and the compiler cannot check it. The program checks it instead,
every time it reads an element, and stops at once when the index is outside the
array.

\begin{lstlisting}[style=bashVerb, escapechar=@+]
@+\textcolor{teal!80}{alice@dev:\mtilde}@+$ ./months
Month: 13
Panic in file "months.yr", at line 10, in function "months::main" !!!
Aborted (core dumped)
\end{lstlisting}

This is a \textit{panic}: the program gives up, and says where. Stopping is the
safe choice. The elements of an array are stored one after the other in memory,
and \token{days[i]} is found by counting \token{i} elements from the first.
Nothing marks the end of the array in memory itself: the element that
\token{days[12]} designates lies right after the last one, where the memory
holds something else, another variable of the program for instance. Without the
check, reading it would give a number that means nothing, and the program would
carry on with it. Writing it would be worse, since it would change that other
variable, which no line of the program names.

A language whose programs cannot touch memory outside what they were given is
said to be \textit{memory safe}. Ymir is, and the check of every index is part
of the price. A program should still check an index that it receives from
outside before using it, to answer with a message rather than a panic;
Section~\ref{sec:collections:ranges} does so for the month.

\notebox{C does not check indices, and writing past the end of an array is the
  \textit{buffer overflow}, one of the most exploited flaws in software. A
  program copies text received from the network into an array of 64
  characters, without checking its length. An attacker sends a longer text,
  whose end overwrites the memory that follows the array, and chooses what is
  written there: among it is the address where the function must go back to
  when it returns. Chapter~\ref{chap:layout} shows what the memory around a
  variable holds, and how the compiler places and aligns the values in it.}

\subsection{Filling an array}

\token{[v ; n]} is an array of \token{n} elements, all equal to \token{v}:
\token{[0 ; 6]} is \token{[0, 0, 0, 0, 0, 0]}. The length must be known while
compiling, since it is part of the type. This form is the usual start of an
array whose elements are computed later, such as the counters of the next
listing, which rolls a die 6\,000 times and counts how often each face comes up.

\lstinputlisting[style=coloredverbatim, caption={Counting the faces of a die, \textit{examples/collections/dice.yr}}, label=lst:collections:dice]{examples/collections/dice.yr}
\vspace{-10pt}%
\begin{lstlisting}[style=bashVerb]
1: 1003
2: 1074
3: 989
4: 945
5: 977
6: 1012
\end{lstlisting}

Each face comes up about 1\,000 times, and the counts change from one run to the
next.

\subsection{Changing the elements}
\label{sec:collections:dmut}

Line~8 of Listing~\ref{lst:collections:dice} adds 1 to one element of
\token{counts}. That is allowed because \token{counts} is declared with
\token{dmut} on line~5. With \token{let} alone, or with \token{let mut}, the
elements of an array cannot change.

\begin{lstlisting}[style=coloredverbatimError, escapechar=@, caption=Changing an element of an array declared with \token{let}]
use std::io;
use std::rand;

fn main() {
  let counts = [0 ; 6];
  for _ in 0 .. 6000 {
    let die = uniform(1, 6);
    @\hb{counts[die - 1]}@ += 1; // error
  }
  println(counts);
}
\end{lstlisting}

The compiler says \errmsg{E4082}{left operand of type i32 is immutable}: the
element, an \token{i32}, is a constant. A collection has two things that can
change: the variable, which can be given a whole new collection, and the
elements inside it. The declaration says which ones may.

\begin{center}
  \begin{adjustbox}{max width=\linewidth}
    \begin{tabular}{|l|l|}
      \hline
      Declaration & What may change\\[0pt]
      \hline
      \hline
      \token{let counts = [0 ; 6];} & nothing\\[0pt]
      \token{let mut counts = [0 ; 6];} & the variable: \token{counts = [1, 1, 1, 1, 1, 1];} is allowed\\[0pt]
      \token{let dmut counts = [0 ; 6];} & the variable and its elements: \token{counts[0] = 5;} is allowed too\\[0pt]
      \hline
    \end{tabular}
  \end{adjustbox}
\end{center}

The \token{d} of \token{dmut} stands for \textit{deep}: the mutability goes all
the way down, to the elements. The same rule applies to the fields of a tuple,
so \token{pair.0 = 5;} needs a \token{pair} declared with \token{dmut}.
Why Ymir tells the two levels apart only shows when several variables refer to
the same elements, which the next chapter is about. Until then, the rule is
enough: a collection whose elements change is declared \token{dmut}.

\subsection{Going through an array}

\token{for x in a} runs its block once for each element of \token{a}, in order
of index, and \token{x} holds the current element. With two names, as on
line~10 of Listing~\ref{lst:collections:dice}, the first holds the index, a
\token{usize}, and the second the element. Line~11 prints \token{i + 1}, so that
the faces are numbered from 1, as on the die.

Both names are constants, like the iterator of a range: \token{count = 0;} in
the block would not reset an element of \token{counts}, and is refused. To
change the elements, the block goes through the indices and writes
\token{counts[i]}, or uses the loops of the next chapter.

\subsection{Expanding an array}
\label{sec:collections:expand_array}

\token{expand} (Section~\ref{sec:collections:expand}) replaces a tuple with its
fields, which requires the compiler to know how many there are. The length of
an array is part of its type, so the compiler knows it too, and an array can be
expanded like a tuple.

\lstinputlisting[style=coloredverbatim, caption={Expanding an array, \textit{examples/collections/seasons.yr}}, label=lst:collections:seasons]{examples/collections/seasons.yr}
\vspace{-10pt}%
\begin{lstlisting}[style=bashVerb]
24
[3, 4, 5, 6, 7, 8], 6 months
(3, 4, 5, spring)
\end{lstlisting}

\begin{itemize}
\item \textit{In a call.} On line~9, \token{volume(expand size)} is the call
  \token{volume(2, 3, 4)}: each element of the array is an argument.
\item \textit{In an array.} On line~13, \token{[expand spring, expand summer]}
  is the array \token{[3, 4, 5, 6, 7, 8]}, the elements of \token{spring}
  followed by those of \token{summer}. Its type is \token{[i32 ; 6]}: the
  compiler adds the lengths.
\item \textit{In a tuple.} On line~16, \token{(expand spring, "spring")} is a
  tuple of four fields, three numbers and a string.
\end{itemize}

The slices of the next section have a length known only when the program runs,
in general, and cannot be expanded whole: a slice received as a parameter is
refused with \errmsg{E4257}{unknown length of expansion for type [i32]}. A part
of a slice whose bounds are known while compiling can be
(Section~\ref{sec:collections:expand_slice}).

\subsection{Arrays of arrays}

The elements of an array can be of any type, arrays included.
\token{[[0 ; 3] ; 3]} is a grid of 3 rows of 3 zeros, of type
\token{[[i32 ; 3] ; 3]}. \token{grid[1]} is its second row, and
\token{grid[1][2]} the third element of that row.

\begin{lstlisting}[style=coloredverbatim]
let dmut grid = [[0 ; 3] ; 3];
grid[1][2] = 5;
println(grid[1]);
for row in grid {
  for cell in row {
    print(cell, " ");
  }
  println();
}
\end{lstlisting}
\vspace{-10pt}%
\begin{lstlisting}[style=bashVerb]
[0, 0, 5]
0 0 0
0 0 5
0 0 0
\end{lstlisting}

\token{grid[1][2] = 5;} changes one element of one row, which needs
\token{grid} declared \token{dmut}, as for any element. The loops go through
the grid in two levels: the outer loop gives each row, an array of 3 elements,
and the inner loop each element of that row.

Exercise~9 plays on such a grid, and Section~\ref{sec:collections:grids} comes
back to grids whose size is known only at run time.

\vfill\pagebreak
\section{A number of values known at run time: slices}
\label{sec:collections:slices}

The length of an array is written in its type, so it must be known while
compiling. A program that keeps the numbers the user types does not know how
many there will be until the user stops. A \textit{slice} is a sequence of
values of the same type whose length is known only when the program runs, and
which can grow.

\lstinputlisting[style=coloredverbatim, caption={Keeping the numbers typed, \textit{examples/collections/numbers.yr}}, label=lst:collections:numbers]{examples/collections/numbers.yr}
\vspace{-10pt}%
\begin{lstlisting}[style=bashVerb, escapechar=@+]
@+\textcolor{teal!80}{alice@dev:\mtilde}@+$ ./numbers
Number (0 to stop): 4
Number (0 to stop): 8
Number (0 to stop): 15
Number (0 to stop): 16
Number (0 to stop): 23
Number (0 to stop): 42
Number (0 to stop): 0
You typed 6 numbers: [4, 8, 15, 16, 23, 42]
\end{lstlisting}

The type of a slice is the type of its elements between brackets, with no
length: \token{numbers} is an \token{[i32]}, on line~4. \token{[]} is the empty
slice, of length 0. Everything else works as for arrays: \token{numbers[i]} is
the element at index \token{i}, counted from 0 and checked when the program
runs, \token{numbers.len} is the length, and \token{for} goes through the
elements, with or without their index.

\subsection{Growing a slice}

The operator \tokennolst{\mtildeascii{}} puts two sequences end to end, and
makes a new slice of their elements, which must have the same type:
\tokennolst{[1, 2] \mtildeascii{} [3]} is \token{[1, 2, 3]}, and neither
operand changes. \tokennolst{numbers \mtildeascii{}= [n]}, on line~10, is
short for \tokennolst{numbers = numbers \mtildeascii{} [n]}: it appends the
elements of the array \token{[n]}, one element here, at the end of
\token{numbers}. The right operand is a sequence, not a single value, so
\tokennolst{numbers \mtildeascii{}= n} is refused.

The operator is not \token{+}, because adding two sequences could as well mean
adding their elements one by one, and \tokennolst{\mtildeascii{}} leaves no
doubt. Since it gives the variable a new value, \tokennolst{\mtildeascii{}=}
needs a variable declared \token{mut}, and \token{dmut} is not needed for it. The
rule of Section~\ref{sec:collections:dmut} is unchanged: \token{numbers[0] = 1;}
would need \token{dmut}.

\subsection{From an array to a slice: \texttt{copy}}
\label{sec:collections:copy}

Square brackets around values always make an array, whatever the type of the
variable that receives it.

\begin{lstlisting}[style=coloredverbatimError, escapechar=@, caption=An array where a slice is expected]
use std::io;

fn main() {
  let odd: [i32] = @\hb{[1, 3, 5, 7]}@; // error
  println(odd);
}
\end{lstlisting}

The compiler says \errmsg{E4095}{incompatible types [i32] and mut [mut i32 ;
  4us]}. The two \token{mut} of the message describe the literal, which the
program may change; the next chapter explains them. What matters is that the
value is an array of four elements, where a slice is expected.

The two are stored differently. The elements of an array are stored in the
variable itself, among the variables of the function that declares it, and
they disappear when that function returns. The elements of a slice are stored
elsewhere, in memory that lasts as long as the slice needs them, and the slice
only records where they are and how many there are. That is why a slice can
grow, and why a function can return one. Making a slice from an array therefore
copies the elements to that other place, and Ymir wants a copy to be written:
\token{copy [1, 3, 5, 7]} is a slice of four elements.

This explanation stays vague on purpose: what \textit{stored in the variable}
and \textit{elsewhere} mean takes knowing how the memory of a program is laid
out. Until Chapter~\ref{chap:layout} gives it, the rule must be accepted as it
is: a slice made from the elements of an array is written with \token{copy}.

\subsection{Slices as parameters}

A parameter of type \token{[i32]} accepts a slice of any length, so one function
serves for all of them. The next listing computes a sum and an average.

\lstinputlisting[style=coloredverbatim, caption={Functions that take a slice, \textit{examples/collections/average.yr}}, label=lst:collections:average]{examples/collections/average.yr}
\vspace{-10pt}%
\begin{lstlisting}[style=bashVerb]
Sum: 98
Average: 14
Weekend average: 22
\end{lstlisting}

An array is not a slice, so it needs a \token{copy} too before it is passed to
\token{average}, on line~21. Without it, the call is refused with
\errmsg{E4226}{the call operator is not defined for \errhole{} and
  \errhole{}}, and the note under it gives the reason,
\errmsg{E4095}{incompatible types [i32] and [i32 ; 2us]}.

Line~12 divides an \token{i32} by a length, which is a \token{usize}. Ymir does
not mix integer types in an operation, and would refuse \token{sum(values) /
  values.len}. \token{cast!i32(x)} converts \token{x} to an \token{i32}: as in
\token{read!i32}, the \token{!i32} names the type wanted. A \token{usize} can hold
numbers too large for an \token{i32}, but the lengths of this book stay far
below two billion, where the conversion would go wrong.

\token{average} has a flaw: the average of an empty slice divides by 0. The
program is then stopped by the system with the message \textit{Floating point
  exception}, although no floating point number is involved.
Section~\ref{sec:collections:options} gives a function a way to say that it has
no answer, and an exercise fixes \token{average} with it.

\subsection{Strings are slices}

Chapter~\ref{chap:types} wrote the type of a string \token{[c8]}. It is indeed a
slice of characters, and everything above applies to it: \token{s.len},
\token{s[0]}, \tokennolst{\mtildeascii{}}, \token{==}, and \token{for c in s}.
\token{read} can read a whole line, when it is asked for a \token{[c8]}. A type
of more than one word goes between braces, so the call is
\token{read!\{[c8]\}}, and returns the line typed, without its line break.

\lstinputlisting[style=coloredverbatim, caption={A string read, joined and gone through, \textit{examples/collections/greeting.yr}}, label=lst:collections:greeting]{examples/collections/greeting.yr}
\vspace{-10pt}%
\begin{lstlisting}[style=bashVerb, escapechar=@+]
@+\textcolor{teal!80}{alice@dev:\mtilde}@+$ ./greeting
Your name: Alice
Hello, Alice! That is 13 characters.
A l i c e
\end{lstlisting}

As Section~\ref{sec:types:encoding} explained, the length of a \token{[c8]} is a
number of bytes, and it is a number of characters only for text in ASCII. A loop
over a \token{[c8]} goes through it byte by byte as well, so a name such as
\textit{Zoé} is not printed as three letters.

\subsection{Part of a slice}

A range between the brackets selects part of a slice. \token{s[a .. b]} is the
slice of the elements from index \token{a} to index \token{b}, \token{b}
excluded, as in the range \token{a .. b}
(Figure~\ref{fig:collections:subslice}). Inside the brackets, \token{\$} stands
for the length of the slice, so \token{s[\$ - 1]} is its last element and
\token{s[2 .. \$]} all of its elements except the first two.

\begin{lstlisting}[style=coloredverbatim]
let primes = copy [2, 3, 5, 7, 11, 13];
println(primes[1 .. 4]); // [3, 5, 7]
println(primes[3 .. $]); // [7, 11, 13]
println(primes[$ - 1]);  // 13
\end{lstlisting}

\begin{figure}[h]
  \centering
  \begin{adjustbox}{max width=\linewidth}
    \begin{tikzpicture}[cell/.style={rectangle, draw=black!60, fill=background!40,
          minimum width=10mm, minimum height=7mm, font=\ttfamily},
        picked/.style={cell, fill=applegreen!25}]
      \foreach \value [count=\i from 0] in {2, 3, 5, 7, 11, 13} {
        \ifnum\i<1
          \node[cell] (c\i) at (\i * 1.0, 0) {\value};
        \else\ifnum\i<4
          \node[picked] (c\i) at (\i * 1.0, 0) {\value};
        \else
          \node[cell] (c\i) at (\i * 1.0, 0) {\value};
        \fi\fi
        \node[cflab, below=1mm of c\i] {\i};
      }
      \node[cflab, left=3mm of c0] {\token{primes}};
      \draw[black!60] ($(c1.north west) + (0.5mm, 1.5mm)$) -- ++(0, 1.5mm)
        coordinate (l) -- ($(c3.north east) + (-0.5mm, 3mm)$) coordinate (r)
        -- ++(0, -1.5mm);
      \node[cflab, above=0.5mm] at ($(l)!0.5!(r)$) {\token{primes[1 .. 4]}};
    \end{tikzpicture}
  \end{adjustbox}
  \caption{\label{fig:collections:subslice} The part \token{primes[1 .. 4]} of a
    slice. It starts at index 1 and stops before index 4, as the range
    \token{1 .. 4} does, so it holds three elements.}
\end{figure}


Taking part of a slice copies nothing: the new slice refers to the same elements
as the first, only fewer of them. The same brackets work on an array. When the
bounds are known while compiling, as in \token{days[0 .. 6]}, the part is an
array of its own, of 6 elements. When they are known only at run time, its
length is unknown, so it can only be a slice, and must be copied,
\token{copy days[0 .. n]}, for the reason given in
Section~\ref{sec:collections:copy}.

\subsection{Expanding part of a slice}
\label{sec:collections:expand_slice}

A slice cannot be expanded whole in general, since its length is known only
when the program runs (Section~\ref{sec:collections:expand_array}). A part of a
slice can, as long as its bounds are known while compiling: \token{s[0 .. 3]}
has 3 elements, whatever the length of \token{s}.

\lstinputlisting[style=coloredverbatim, caption={Expanding part of a slice, \textit{examples/collections/boxes.yr}}, label=lst:collections:boxes]{examples/collections/boxes.yr}
\vspace{-10pt}%
\begin{lstlisting}[style=bashVerb, escapechar=@+]
@+\textcolor{teal!80}{alice@dev:\mtilde}@+$ ./boxes
24
120
Panic in file "boxes.yr", at line 8, in function "boxes::firstVolume" !!!
Aborted (core dumped)
\end{lstlisting}

On line~8, \token{expand sizes[0 .. 3]} passes the first three elements of
\token{sizes} as the three arguments of \token{volume}. \token{sizes} is a
parameter, of any length, but the part has 3 elements, and the compiler knows
it. On line~14, the function receives the slice \token{[4, 5, 6]}, whose first
three elements give 120.

On line~15, the slice has only 2 elements, and the part \token{sizes[0 .. 3]}
does not exist. As for an index, the compiler cannot check it while compiling,
and the program checks it when it runs: it panics.

The bounds must be known while compiling. With bounds that are variables, as in
\token{expand sizes[n .. m]}, the length of the part is known only when the
program runs, and the compiler refuses it with \errmsg{E4257}{unknown length of
  expansion for type mut [i32]}.

\subsection{Slices of slices}
\label{sec:collections:grids}

The elements of a slice can be slices. A grid whose size is known only at run
time, such as a table of \token{rows} rows and \token{cols} columns, can be a
slice of \token{rows} slices of \token{cols} elements each, of type
\token{[[i32]]}. As for an array of arrays, \token{grid[i]} is row \token{i},
and \token{grid[i][j]} its element in column \token{j}.

\begin{lstlisting}[style=coloredverbatim]
let rows = 3;
let cols = 4;
let dmut grid: [[i32]] = [];
for _ in 0 .. rows {
  grid ~= [copy [0 ; cols]];
}
grid[1][2] = 7;
println(grid); // [[0, 0, 0, 0], [0, 0, 7, 0], [0, 0, 0, 0]]
\end{lstlisting}

Each row is a slice of its own, made by a \token{copy}, and its elements are
stored away from the variable (Section~\ref{sec:collections:copy}). Reserving
such a block of memory is called an \textit{allocation}. The grid above takes
one allocation per row, and more for \token{grid} itself, which only records
where each row is and how long it is. The rows are not next to each other in
memory: each lies wherever its block was found
(Figure~\ref{fig:collections:grid_layout}, above).

\cautionbox{\token{copy [copy [0 ; cols] ; rows]} looks like a shorter way to make
  the grid, but it does not make \token{rows} rows. \token{[v ; n]} computes
  \token{v} once, and repeats it \token{n} times: every element of the grid is
  then the same row, and \token{grid[1][2] = 7;} changes the third element of
  every row. The loop above makes a new row at each turn. The next chapter
  explains why, and how \token{dcopy} makes the rows distinct.}

\subsection{A grid in a single slice}
\label{sec:collections:flat_grid}

An allocation takes time, and a grid of a thousand rows takes a thousand of
them. A single slice of \token{rows * cols} elements holds the same numbers
with one allocation: \token{copy [0 ; rows * cols]}. Since the result of
\token{copy} is a slice, whose length is not part of its type, the \token{n} of
\token{[v ; n]} may then be known only at run time.

The program lays the rows out one after the other, and finds the elements
itself. Row 0 takes the indices 0 to \token{cols - 1}, row 1 the next
\token{cols}, and so on. The element at row \token{i} and column \token{j}
comes after \token{i} full rows of \token{cols} elements, and after \token{j}
elements of its own row: it is at index \token{i * cols + j}
(Figure~\ref{fig:collections:grid_layout}, below).

\begin{lstlisting}[style=coloredverbatim]
let rows = 3;
let cols = 4;
let dmut flat = copy [0 ; rows * cols];
flat[1 * cols + 2] = 7;
println(flat); // [0, 0, 0, 0, 0, 0, 7, 0, 0, 0, 0, 0]
\end{lstlisting}

\begin{figure}[h]
  \centering
  \begin{adjustbox}{max width=\linewidth}
    \begin{tikzpicture}[cell/.style={rectangle, draw=black!60, fill=background!40,
          minimum width=10mm, minimum height=7mm, font=\ttfamily},
        ref/.style={cell, minimum width=14mm, font=\footnotesize},
        picked/.style={cell, fill=applegreen!25}]
      \node[ref] (o0) at (0, 0) {row 0};
      \node[ref] (o1) at (0, -0.7) {row 1};
      \node[cflab, above=1mm of o0] {\token{grid}};

      \foreach \value [count=\j from 0] in {1, 2, 3} {
        \node[cell] (a\j) at (2.5 + \j * 1.0, 0.6) {\value};
      }
      \foreach \value [count=\j from 0] in {4, 5, 6} {
        \ifnum\j=2
          \node[picked] (b\j) at (4.5 + \j * 1.0, -1.3) {\value};
        \else
          \node[cell] (b\j) at (4.5 + \j * 1.0, -1.3) {\value};
        \fi
      }
      \draw[cfflow] (o0.east) -- (a0.west);
      \draw[cfflow] (o1.east) -- (b0.west);
      \node[cflab, right=3mm of b2] {\token{grid[1][2]}};

      \foreach \value [count=\i from 0] in {1, 2, 3, 4, 5, 6} {
        \ifnum\i=5
          \node[picked] (f\i) at (\i * 1.0 + 1.0, -3.4) {\value};
        \else
          \node[cell] (f\i) at (\i * 1.0 + 1.0, -3.4) {\value};
        \fi
        \node[cflab, below=1mm of f\i] {\i};
      }
      \node[cflab, left=3mm of f0] {\token{flat}};
      \draw[black!60] ($(f0.north west) + (0.5mm, 1.5mm)$) -- ++(0, 1.5mm)
        coordinate (l0) -- ($(f2.north east) + (-0.5mm, 3mm)$) coordinate (r0)
        -- ++(0, -1.5mm);
      \node[cflab, above=0.5mm] at ($(l0)!0.5!(r0)$) {row 0};
      \draw[black!60] ($(f3.north west) + (0.5mm, 1.5mm)$) -- ++(0, 1.5mm)
        coordinate (l1) -- ($(f5.north east) + (-0.5mm, 3mm)$) coordinate (r1)
        -- ++(0, -1.5mm);
      \node[cflab, above=0.5mm] at ($(l1)!0.5!(r1)$) {row 1};
      \node[cflab, right=3mm of f5] {\token{flat[1 * 3 + 2]}};
    \end{tikzpicture}
  \end{adjustbox}
  \caption{\label{fig:collections:grid_layout} Two ways to keep a grid of 2 rows
    of 3 elements. Above, a slice of slices: \token{grid} records where each row
    is, and each row is a block of memory of its own, wherever it was
    allocated. Below, a single slice: the rows lie one after the other in one
    block, and the element at row 1 and column 2 comes after one full row of 3
    elements and 2 elements of its own row, at index $1 \times 3 + 2 = 5$.}
\end{figure}


The next listing prints a multiplication table of the size the user asks for,
kept in a single slice.

\lstinputlisting[style=coloredverbatim, caption={A grid in a single slice, \textit{examples/collections/table.yr}}, label=lst:collections:table]{examples/collections/table.yr}
\vspace{-10pt}%
\begin{lstlisting}[style=bashVerb, escapechar=@+]
@+\textcolor{teal!80}{alice@dev:\mtilde}@+$ ./table
Rows: 4
Columns: 5
1 2 3 4 5
2 4 6 8 10
3 6 9 12 15
4 8 12 16 20
Row 3, column 4: 12
\end{lstlisting}

The loops of lines~8 to~12 write the element of row \token{i} and column
\token{j} at index \token{i * cols + j}, and the loops of lines~14 to~19 read it
back from there. On line~20, the element of the third row and the fourth
column, counted from 1 by the user, is at row 2 and column 3, counted from 0:
index \token{2 * cols + 3}, 13 here.

\vfill\pagebreak
\section{Ranges are values}
\label{sec:collections:ranges}

The ranges of Section~\ref{sec:control:for} were written in \token{for} loops,
and in the brackets of the previous section. They are values in their own right,
which a variable can hold and a function can take or return. The type of a range
of \token{i32} is \token{..i32}.

\begin{lstlisting}[style=coloredverbatim, caption=A range stored in a variable, label=lst:collections:range_value]
let week = 1 ... 7;
println(week.fst, " to ", week.scd); // 1 to 7
println(3 in week, " ", 9 in week);  // true false

let primes = copy [2, 3, 5, 7, 11, 13];
let middle = 1 .. 4;
println(primes[middle]);             // [3, 5, 7]
\end{lstlisting}

\token{.fst} and \token{.scd} are the first and the second bound of the range,
and \token{.contains} says whether the second one belongs to it: it is
\token{false} for a range written with \token{..}, and \token{true} for a range
written with \token{...}, such as \token{week}.

\subsection{Testing that a value is in a range}

\token{x in r} is \token{true} when the value \token{x} is one of those that the
range \token{r} goes through, and \token{x !in r} is its opposite. It says in
one expression what \token{x >= 1 \&\& x <= 7} says in two comparisons, with the
bounds written only once. It is the check that the months of
Listing~\ref{lst:collections:months} were missing.

\begin{lstlisting}[style=coloredverbatim, escapechar=@, caption=Checking a month before using it as an index, label=lst:collections:months_checked]
use std::io;

fn main() {
  let names = ["January", "February", "March", "April",
               "May", "June", "July", "August",
               "September", "October", "November", "December"];
  let days = [31, 28, 31, 30, 31, 30, 31, 31, 30, 31, 30, 31];

  let month = read!i32(ask-> "Month: ");
  @\hcb{if month !in 1 ... 12 \{}@
    @\hcb{println("There is no month ", month, ".");}@
  @\hcb{\} else \{}@
    println(names[month - 1], " has ", days[month - 1], " days.");
  @\hcb{\}}@
}
\end{lstlisting}
\vspace{-10pt}%
\begin{lstlisting}[style=bashVerb, escapechar=@+]
@+\textcolor{teal!80}{alice@dev:\mtilde}@+$ ./months
Month: 13
There is no month 13.
\end{lstlisting}

\subsection{Testing that a value is in a sequence}

\token{in} also tests whether a value is among the elements of an array or of a
slice, and \token{!in} whether it is not.

\begin{lstlisting}[style=coloredverbatim]
let primes = copy [2, 3, 5, 7, 11, 13];
println(7 in primes, " ", 8 in primes);  // true false
println('y' in "Ymir", " ", 'Y' in "Ymir");  // false true
\end{lstlisting}

The test goes through the elements one by one, until it finds the value or
reaches the end. It says whether the value is there, not where: finding the
index of a value takes a loop, which the next section writes.

\vfill\pagebreak
\section{A value that may be missing: options}
\label{sec:collections:options}

Searching a slice for a value has two outcomes: the value is at some index, or
it is not there at all. A function that returns the index as a \token{usize}
has nothing to return in the second case. Some languages return a special
number, such as $-1$, and count on the caller to test it. A \token{usize} cannot
be negative, and a caller who forgets the test uses $-1$ as an index. An
\textit{option} is a value that holds either a value of a given type, or nothing,
and the compiler makes sure that the program tells the two apart.

\lstinputlisting[style=coloredverbatim, caption={Searching a slice, \textit{examples/collections/find.yr}}, label=lst:collections:find]{examples/collections/find.yr}
\vspace{-10pt}%
\begin{lstlisting}[style=bashVerb, escapechar=@+]
@+\textcolor{teal!80}{alice@dev:\mtilde}@+$ ./find
Ok(3) Err()
Number: 11
11 is prime number 5.
@+\textcolor{teal!80}{alice@dev:\mtilde}@+$ ./find
Ok(3) Err()
Number: 12
12 is not among the first eight primes.
\end{lstlisting}

The type of an option is the type of its value followed by a question mark.
On line~3, \token{find} returns a \token{usize?}, an option that may hold a
\token{usize}.

An option that holds a value is written the same way, with the question mark
after the value. On line~5, \token{i?} is the option that holds the index
\token{i}, and \token{find} returns it as soon as the value is found.

The option that holds nothing is written with the keyword \token{none}, on
line~7. It takes its type from where it goes, here the result of
\token{find}.

\token{println} prints an option that holds a value as \token{Ok(3)}, and an
option that holds nothing as \token{Err()}.

\subsection{A loop whose block is an \texttt{if}}

Line~4 writes the \token{for} and the \token{if} on the same line, with a
single block. The block of a loop may be replaced by one expression, and an
\token{if} is one: the line reads \textit{for each \token{v} of
\token{values}, if \token{v} is \token{wanted}}. When the whole block of a
loop is an \token{if} without \token{else}, Ymir programs write it this way,
which saves a level of indentation. An \token{if} written this way has no
\token{else}: when one is needed, the \token{if} goes inside the braces of the
loop, as in the previous chapters.

\subsection{Getting the value out}

An option is not the value it may hold, and cannot be used as one.

\begin{lstlisting}[style=coloredverbatimError, escapechar=@, caption=An option used as a number]
use std::io;

fn find(values: [i32], wanted: i32)-> usize? {
  for i, v in values if v == wanted {
    return i?;
  }
  none
}

fn main() {
  let primes = copy [2, 3, 5, 7, 11, 13, 17, 19];
  let i = find(primes, 7);
  println("7 is prime number ", @\hb{i + 1}@); // error
}
\end{lstlisting}

The compiler says \errmsg{E4224}{undefined operator + for types (usize)? and
  i32}. That is the whole point of an option: the value can be reached only
through code that also says what to do when it is missing. The usual way is a
\token{match}, as on lines~15 to~18 of Listing~\ref{lst:collections:find}. The
pattern \token{Ok(i)} fits an option that holds a value, and names that value
\token{i} for the arm; \token{Err()} fits an option that holds nothing.

When only one of the two cases needs work, \token{if let} is shorter. It tries a
single pattern, and runs its block when the pattern fits, with the names the
pattern declares; the \token{else} block, if any, runs otherwise.

%% check: skip -- partial program, find is the function of the listing above
\begin{lstlisting}[style=coloredverbatim]
if let Ok(i) = find(primes, n) {
  println(n, " is prime number ", i + 1, ".");
} else {
  println(n, " is not among the first eight primes.");
}
\end{lstlisting}

\token{o.hasValue} is \token{true} when the option \token{o} holds a value, for
a test that does not need the value itself.

An option also has a field \token{o.value}, which reads the value directly. When
the option holds nothing, it does not panic but throws an \textit{exception},
which a later chapter of the course presents, and Part II in
Chapter~\ref{chap:error}. Until then, \token{match} and \token{if let} do the
same work safely.

\notebox{An empty option is printed \token{Err()} rather than \token{None()},
  because an option can do more than hold nothing: it can hold the reason why it
  has no value. \token{find} has only one reason to fail, so it gives none, but a
  function that reads a file can fail because the file does not exist, or
  because it may not be read. That use of options belongs to error handling.}

